\chapter{Canais e produtos}

\section{Canais oficiais}

Os canais oficiais de representação do Conselho Regional de Psicologia
de São Paulo (CRP~SP) são os seguintes:

\begin{enumerate}

  \item
        \emph{site} e \emph{e-mail} institucional;
  \item
        perfis oficiais em redes sociais;
  \item
        eventos promovidos pela autarquia, em parceria ou não, nos estados e
        territórios.
\end{enumerate}

\subsection{Site institucional}

O \emph{site} institucional é o principal repositório de informações do
CRP~SP. Nele devem ser disponibilizadas informações sobre todas as
atividades do Conselho, além de material de referência e orientação para
a categoria. O \emph{site} também deve oferecer canais para contato com
todas as unidades da autarquia.

\subsection{Redes sociais}

O CRP~SP mantém perfis ativos nas seguintes redes sociais:

\begin{enumerate}

  \item
        Facebook;
  \item
        Instagram;
  \item
        LinkedIn;
  \item
        X/Twitter;
  \item
        YouTube.
\end{enumerate}

Os perfis sociais do CRP~SP têm como objetivo formar e informar a
categoria sobre temas associados à Psicologia, à atuação profissional e
ao funcionamento do Sistema Conselhos. Por meio deles, buscamos ampliar
nossos pontos de contato e comunicação com psicólogas de todo o estado
de São Paulo, além de participar dos debates sobre a Psicologia como
ciência e profissão.

Em observância aos princípios da transparência e a nosso compromisso
ético, estabelecemos as seguintes diretrizes:

\begin{enumerate}

  \item
        \textbf{perfis e páginas oficiais}: só são permitidos perfis e páginas
        associados ao CRP~SP que sejam criados e mantidos pela instituição. A
        autarquia se reserva o direito de não endossar perfis e páginas que se
        autointitulem associados ao Conselho e que não sejam reconhecidos como
        oficiais, tomando as devidas medidas quando necessário;
  \item
        \textbf{gestão das contas}: as contas oficiais do CRP~SP serão
        gerenciadas por uma equipe certificada e autorizada, composta por
        profissionais de comunicação, trabalhadoras e conselheiras da
        autarquia;
  \item
        \textbf{política de interação e moderação}: para garantir a
        veracidade, exatidão, confiabilidade, qualidade e legalidade das
        informações, as interações nos canais oficiais estarão sujeitas a
        moderação. Mensagens e perfis poderão ser excluídos se:

        \begin{enumerate}

          \item
                usarem linguagem inapropriada, obscena, caluniosa, grosseira,
                abusiva, difamatória ou ofensiva;
          \item
                fizerem apologia a práticas ilícitas;
          \item
                incitarem ódio, violência, racismo ou qualquer forma de
                discriminação;
          \item
                contiverem ameaças, assédio, injúria, calúnia, difamação ou qualquer
                outro ilícito penal;
          \item
                divulgarem conteúdos na forma de \emph{spam} ou ``correntes'';
          \item
                tiverem objetivo comercial ou publicitário;
          \item
                forem ininteligíveis ou fora de contexto;
          \item
                contiverem propagandas político-partidárias;
          \item
                contiverem \emph{links} suspeitos ou representarem ameaça à
                segurança da informação;
          \item
                usarem informações e imagens de pessoas e instituições de modo
                indevido;
          \item
                contiverem dados pessoais da autora, do autor ou de terceiros;
          \item
                violarem direitos de imagem e de propriedade intelectual;
          \item
                forem fraudulentos (perfis falsos e/ou \emph{bots}) ou promoverem
                conteúdo inverídico (\emph{fake news}).
        \end{enumerate}
  \item
        \textbf{restrições de divulgação}: o CRP~SP não segue e não marca
        pessoas, veículos de comunicação e empresas privadas: apenas entidades
        parceiras e órgãos do poder público. Não divulgamos partidos políticos
        e não realizamos atendimento administrativo e de orientação nas redes
        sociais. As demandas que não se enquadrarem nos produtos fixos devem
        ser levadas a comissões estaduais.
\end{enumerate}

A não observância dessas diretrizes pode acarretar ações imediatas do
CRP~SP, sem necessidade de justificativa, consulta ou aviso prévio.
Mensagens inadequadas poderão ser encaminhadas às autoridades
responsáveis. Nosso compromisso é com a ética, a transparência e a
qualidade das nossas comunicações.

\section{Produtos}

Dividimos os produtos de comunicação do CRP~SP em \emph{produtos de
  circulação rápida} e \emph{produtos de circulação lenta} (ou de ``cauda
longa''). Entre os primeiros, incluem-se conteúdo produzido para redes
sociais (como o CRP~SP Acontece e o CRP~SP na Mídia); os de circulação
lenta são aqueles com potencial para se tornar material de referência e
consulta (Cadernos Temáticos, Jornal Psi, vídeos institucionais).

A descrição de produtos de comunicação presente nesta \textbf{Política}
não impede a elaboração de diretrizes e políticas editoriais específicas
e completas para cada produto.

\subsection{Psicologia em Ação}

Divulgação de eventos a pedido, utilizando-se do \emph{site} (Agenda
CRP~SP), de Stories (perfil oficial no Instagram), postagens (demais
redes sociais) e \emph{cards} a serem compartilhados pelos territórios.
A cobertura dos eventos é publicada no CRP~SP Acontece.

As publicações de \emph{Psicologia em Ação} possuem identidade visual
própria, que pode receber elementos temáticos.

\subsection{CRP~SP Acontece}

Veiculado semanal ou quinzenalmente, informa a categoria e a sociedade
sobre as ações, eventos, representações, mobilizações, novidades etc.
promovidos nos territórios, pelas comissões gestoras, em formato de
galeria de imagens (até 10 fotos por publicação), acompanhada de breves
relatos.

Subsedes, comissões e demais instâncias devem enviar seus relatórios
preenchidos e acompanhados de foto(s) para a Comunicação.

\subsection{CRP~SP Informa}

Ferramenta de \emph{e-mail marketing} destinada a dar publicidade às
ações realizadas nos territórios, por meio de envio de mensagens
institucionais a psicólogas e entidades afins. É uma importante
ferramenta de comunicação direcionada, com média estadual de abertura
dos \emph{e-mails} de 20\%.

A cada território é facultado o envio de até quatro mensagens por mês.

\subsection{CRP~SP Articula}

Traz informações sobre parcerias do CRP~SP com outras instituições,
estatais ou de interesse público.

\subsection{CRP~SP Representa}

Divulgação da participação do CRP~SP em eventos para o qual a autarquia
foi convidada oficialmente.

\subsection{CRP~SP Debate (lives)}

Canal que possibilita levar as discussões para o público amplo e trazer
convidadas de diferentes espaços. As \emph{lives} são realizadas por
comissões e, geralmente, estão inseridas em uma campanha. Contam sempre
com tradução para Libras e têm o \emph{chat} mediado pela Coordenação de
Comunicação. São operadas no StreamYard (transmissão pelo YouTube e
Facebook), podendo apresentar até 7 participantes simultaneamente.

\emph{Lives} devem ser solicitadas com no mínimo 30 dias de
antecedência, por meio do formulário de eventos, que deve ser preenchido
e enviado para as pessoas responsáveis por eventos na Coordenação de
Comunicação. Após o recebimento, será criado grupo de imMail para
realização de testes prévios, elaboração de roteiro e alinhamento
necessário.

\subsection{Calendário de datas
  comemorativas}

Artigos e postagens publicados em datas simbólicas para a defesa dos
Direitos Humanos. Textos são organizados pelas comissões permanentes e
encaminhados à ComCom com até três dias úteis de antecedência para
publicação, e devem ser elaborados com linguagem inclusiva
(\#paratodosverem), fornecendo informações sobre o significado da data
em destaque, abordando aspectos relevantes do exercício profissional da
Psicologia e fazendo referência à(s) normativa(s) do Sistema Conselhos
e/ou do Código de Ética alinhada(s) ao assunto da data. A escrita deve
ser objetiva, simples e compatível com as redes sociais (máximo de 1.800
caracteres com espaços e quatro parágrafos). O conteúdo é padronizado e
institucional.

\subsection{CRP~SP na Mídia}

Divulga as participações do CRP~SP na imprensa e em outros canais
externos.

Solicitações de divulgação devem ser encaminhadas para o \emph{e-mail}
\href{mailto:ascom@crpsp.org.br}{ascom@crpsp.org.br}. Os pedidos
são levados para representações e plenário, para que se decida sobre o
encaminhamento.

\subsection{Comunicação integrada}

A comunicação das demais unidades administrativas do Conselho com a
categoria e com outras instituições pode ser adequada à linguagem
institucional mediante solicitação à unidade de Comunicação, observado o
cronograma e o planejamento de atividades desta unidade. A solicitação
deverá ser feita, por \emph{e-mail}, pela pessoa coordenadora da unidade
administrativa solicitante.

\subsection{Notas de pesar}

As notas de pesar têm como objetivo prestar solidariedade e
reconhecimento à pessoa falecida, especialmente quando esta desempenhou
papel significativo para a Psicologia, a defesa dos Direitos Humanos ou
outras causas alinhadas às diretrizes institucionais do CRP~SP.

A redação deve ser respeitosa, empática e compatível com a comunicação
institucional do Conselho, destacando, de forma breve, as contribuições
da pessoa homenageada. O texto não deve conter cunho político-partidário
ou religioso, garantindo imparcialidade e adequação a um público
diversificado.

\subsubsection{Solicitação e
  aprovação}

As solicitações para publicação devem ser encaminhadas à Unidade de
Comunicação, acompanhadas de informações completas, como nome, data de
falecimento e breve descrição das contribuições da pessoa homenageada.

A aprovação do texto será realizada pela Comissão de Comunicação
Institucional (ComCom) ou, quando necessário, pelo Plenário.

\subsubsection{Canais de divulgação}

As notas de pesar serão divulgadas nos canais institucionais do CRP~SP,
conforme a relevância do caso e deliberação da ComCom.

Publicações de notas de pesar devem ser realizadas em caráter
excepcional, preservando a sensibilidade do público e evitando
sobrecarga nos meios de comunicação do Conselho.

\subsubsection{Prazos e responsabilidades}

A Unidade de Comunicação se compromete a publicar a nota em até
\textbf{24 horas (em dias úteis)} entre a solicitação e a validação
final.

O envio da solicitação deve ser feito com antecedência, sempre que
possível, para viabilizar o planejamento adequado.

\subsection{Informe interno}

Mensalmente, a Unidade de Comunicação publica um informe interno, com
temas relevantes para as trabalhadoras do Conselho. Textos e notícias
podem ser enviados para
\href{mailto:ascom@crpsp.org.br}{\nolinkurl{ascom@crpsp.org.br}} com o
assunto ``Informe interno'', para que seja avaliada a viabilidade de sua
publicação.

\subsection{Publicações diárias}

A Coordenação de Comunicação do CRP~SP é responsável pela publicação, a
qualquer tempo e de acordo com a demanda, de:

\begin{itemize}
  \item
        material do Sistema Conselhos de Psicologia;
  \item
        orientações do Centro de Referências Técnicas em Psicologia e
        Políticas Públicas (Crepop);
  \item
        eventos e outros assuntos de interesse produzidos por entidades
        parceiras;
  \item
        notas de posicionamento;
  \item
        informes administrativos e financeiros;
  \item
        orientações gerais.
\end{itemize}

\subsection{Jornal Psi}

O Jornal Psi é publicado em quatro edições anuais no formato
\emph{on-line}, com envio de versões impressas para psicólogas e
entidades com inscrição ativa no CRP~SP.

As pautas são definidas em reuniões entre a Coordenação de Comunicação e
a ComCom. Nessas ocasiões, também é definida a abordagem para cada tema.

A validação da versão final do jornal é feita pelo Plenário.

\subsubsection{Linha editorial}

Atualmente, as seções do jornal incluem:

\begin{itemize}
  \item
        ``Um dia na vida'';
  \item
        ``Cotidiano'';
  \item
        artigo sobre tema pertinente à edição;
  \item
        ``Perspectiva da/o usuária/o'';
  \item
        ``Ética'';
  \item
        ``Orientação'';
  \item
        ``Penalidades éticas'';
  \item
        ``Estante''.
\end{itemize}

\paragraph{Distribuição}

O \emph{mailing} da categoria (nomes/nomes sociais e endereços) para o
envio do Jornal Psi é gerado pelo sistema BRC, a base de dados do
Conselho. A tiragem total excede os 160 mil exemplares. Recebem o Jornal
Psi:

\begin{itemize}
  \item
        psicólogas e empresas com inscrições ativas no CRP~SP;
  \item
        instituições de ensino;
  \item
        Sistema Conselhos de Psicologia;
  \item
        bibliotecas, instituições e órgãos alinhados com a temática da
        Psicologia.
\end{itemize}

\subsection{Identidade visual da
  gestão}

A Coordenação de Comunicação pode, mediante solicitação, incorporar o
mote da gestão da autarquia à identidade visual do CRP~SP.

Com o objetivo de evitar ambiguidades na comunicação, e em observância
ao princípio da impessoalidade, a partir de 2025 não produziremos
logomarcas, assinaturas gráficas ou quaisquer produtos visuais
específicos para a gestão.

\subsection{Cadernos Temáticos}

\textbf{Cadernos Temáticos} é uma série de publicações que trazem o
conteúdo de debates e eventos temáticos promovidos pelo CRP~SP. São
produzidos de acordo com a demanda das comissões permanentes e
especiais, às quais cabe a produção do conteúdo bruto e a organização
lógica do material (separação em capítulos, divisão e identificação das
falas transcritas etc.). As comissões devem garantir a qualidade e a
consistência do texto antes de encaminhar o documento à Unidade de
Comunicação.

A Unidade de Comunicação é responsável pela revisão textual, diagramação
e produção gráfica dos \textbf{Cadernos}. Não cabem à Unidade a correção
do conteúdo bruto, a revisão de transcrições e a produção de textos
originais.

\subsubsection{Atribuições}

\paragraph{Unidade de
  Comunicação}

\begin{enumerate}

  \item
        \textbf{Revisão gramatical e adequação de gênero}:~realizar a revisão
        gramatical e estilística do material consolidado, garantindo a
        adequação de gênero e linguagem inclusiva;
  \item
        \textbf{diagramação}:~desenvolver a diagramação do conteúdo, seguindo
        as diretrizes visuais da instituição;
  \item
        \textbf{produção final}:~preparar os arquivos finais para impressão ou
        publicação digital, conforme a demanda da comissão responsável;
  \item
        \textbf{execução técnica}:~conduzir ajustes técnicos necessários (como
        diagramação de textos aprovados) após o recebimento do material
        validado pela comissão.
\end{enumerate}

\paragraph{Comissões solicitantes}

\begin{enumerate}

  \item
        \textbf{Produção do conteúdo:}~elaborar, consolidar e revisar
        internamente o material bruto antes do envio à unidade de comunicação.
        É responsabilidade da comissão entregar o material em formato adequado
        para a revisão técnica e diagramação;
  \item
        \textbf{organização e compilação:}~organizar os textos, transcrições
        (removendo inconsistências, pausas verbais e interjeições) e demais
        documentos necessários;
  \item
        \textbf{validação:}~aprovar o conteúdo final revisado e diagramado
        antes da publicação ou impressão.
\end{enumerate}

\subsubsection{Estrutura}

A estrutura dos \textbf{Cadernos Temáticos} deve necessariamente conter:

\begin{enumerate}

  \item
        apresentação da \textbf{coleção} Cadernos Temáticos, assinada pela
        diretoria ou pelo Plenário (até 2.500 caracteres, com espaços);
  \item
        \textbf{apresentação} \textbf{da edição}, assinada pela diretoria,
        pelo Plenário ou por representante de comissão relacionada ao tema
        (até 2.000 caracteres, com espaços);
  \item
        \textbf{``miolo''} --- os capítulos propriamente ditos, na forma de
        transcrições de evento, artigos, entrevistas e referências técnicas.
\end{enumerate}

Além disso, cada capítulo deve necessariamente conter:

\begin{enumerate}

  \item
        título;
  \item
        nome da/o autora/or ou palestrante, seguido de minicurrículo;
  \item
        transcrição da fala de cada palestrante ou texto do artigo;
  \item
        transcrição do debate (quando se tratar de publicação de evento), com
        identificação da pessoa que toma a palavra, sempre que possível.
\end{enumerate}

Na publicação de eventos transcritos, o mesmo título pode se referir a
diferentes falas, como acontece em mesas-redondas. Nesse caso, antes da
primeira intervenção de cada palestrante devem ser indicados nome e
minicurrículo. Nas falas subsequentes, basta indicar o nome.

Nos Cadernos Temáticos, também é comum o emprego de ``olhos'' (citações
de partes do texto em destaque).

\subsubsection{Prazo}

Uma vez recebido o texto dos \emph{Cadernos} para publicação, a
Coordenação de Comunicação deve entregar o produto finalizado (impresso
ou não) em no mínimo 30 e no máximo 60 dias.